\section{Final Verdict}
\label{sec:final-verdict}

The initial question asked was whether extensible, multi-source and explainable screening systems could support fairer, more reliable technical
screening compared to single-source rigid black-box alternatives (including typical commercial ATS rankings).
The motivation was to move beyond CV-only screening, with its parsing fragility, self-report bias and opaque automated rankings, and instead
pivot towards demonstrated evidence and recruiter-visible scoring.

We believe that MERIT has achieved these objectives as a working prototype, with a strong foundation for further development. Reflecting on the
results of our evaluation, we see that MERIT is a feasible research-prototype. However, we also acknowledge that it is currently not ready for 
deployment. 

To anchor this judgement, consider the high-discrimination cohort of Study~07 ($n=10$), where MERIT achieved the highest observed expert rank agreement
($\rho = 0.733$, $p = 0.016$) and recovered candidates whom CV-only baselines discarded when GitHub demonstrated role-relevant work (Jordan Smith, Study~01).
However, we should still note that under deliberate CV and GitHub dissonance, MERIT goes strongly against the expert consensus (Study~10).
While MERIT successfully changes who is surfaced under favourable conditions, it does not guarantee a correct ranking in all scenarios.
The main value of the system is that it gives recruiters evidence they can inspect and challenge, rather than a single score to accept without question.

MERIT is intended to be used as an auditable screening assistant. Recruiters are expected to inspect the evidence provided by the system, treat the 
ranks as hypotheses, and not as the final decision. Interpreting MERIT as an oracle is not recommended, not only because of the potential for automation
bias, but also due to it being close to violating Article~22 of the GDPR. To reiterate, this article states that individuals should not be subject to a decision
made entirely by an algorithm without human intervention \cite{eu_gdpr_2016}.

Regarding the three pillars, extensibility and explainability are implemented and partially validated. The multi-source fusion pillar has been shown 
to be a strong differentiator where public evidence is available, but it introduces open-source visibility bias when candidates lack public evidence.
Ultimately, the system's real-world effectiveness remains heavily dependent on both the ingestion's scale and quality.

We validated the feasibility of the prototype by using skill-extraction $F_1$, Spearman~$\rho$, and usability scores as indicators.
However, we must note that these scores cannot serve as proof of improved hire quality, as they are not outcome-linked.

Taken together, MERIT successfully meets its core objective of designing and validating a multi-source screening architecture.
While the current prototype is not yet ready for real-world deployment, it establishes a strong foundation for multi-source, explainable screening.
By addressing the technical and ethical constraints identified in this research, future work can pivot MERIT from a research prototype into a robust, production-ready system.
